% ---------
% --------
%!TEX TS-program = pdflatex
%!TEX encoding = UTF-8 Unicode
%!TEX spellcheck = English (Aspell)

\documentclass[11pt]{article}
\usepackage{jeffe,handout,graphicx}
\usepackage{fourier-orns}
\usepackage{mathtools}
\usepackage[normalem]{ulem} % either use this (simple) or
\usepackage{tikz}

\def\To{\leadsto}
\def\Tostar{\mathrel{\To\!\!\!^*}}
\def\derives{\Tostar}

\def\Cdot{\mathbin{\text{\normalfont \textbullet}}}
\def\Sym#1{\texttt{\color{BrickRed}#1}}
\def\BlueSym#1{\textbf{\texttt{\color{RoyalBlue}#1}}}
\allowdisplaybreaks

\begin{document}

\begin{center}
\Large\textbf{CSCI 432/532, Spring 2026}%
\\
\LARGE\textbf{Homework 5}%
\\[0.5ex]
\large Due Monday, February 23 at 8:30am
\end{center}

\bigskip
\hrule
\bigskip

\subsection*{Submission Requirements}
\begin{itemize}
    \item Type or clearly hand-write your solutions into a PDF format so that they are legible and professional. Submit your PDF on Canvas. 
	\item Put each \textbf{sub}problem on a separate page and select the corresponding problem via the Gradescope interface.
    \item Do not submit your first draft. Type or clearly re-write your solutions for your final submission. If your submission is not legible, we will ask you to resubmit.
\end{itemize}

\subsection*{Collaboration and Resources}

The \emph{best} resources for completing the homework are (in no particular order):
\begin{itemize}
	\item Your own notes from lectures and problem sessions, and/or lecture recordings.
	\item The lecture notes linked from the course website for the lecture day.
	\item The questions channel on Discord.
	\item Office hours.
	\item Discussions with classmates.
\end{itemize}
You may also access almost any resource in preparing your solution to the
homework. The exception is that you may not seek answers to the problems on the
homework directly, such as by pasting them into a GenAI tool, searching for
solutions on a search engine, looking for them on Chegg or similar websites, or
asking another person for the solution.

Additionally, you \EMPH{must}
\begin{itemize}
    \item Write your solutions in your own words---this means that you should never be \emph{copying}
    anything of substance from any source, and
    \item credit every resource you use, including GenAI tools, by providing links or full chat transcripts.
    If you use the provided LaTeX template, you can use the \texttt{Sources}
    environment for this. Ask if you need help!
\end{itemize}


\begin{boxquote}{0.9}
Remember, if you choose to use GenAI tools, not only must you provide a full
transcript of your conversation (or a link to one), you must also use the
following prompt (or some variation of it), and provide a full transcript of
the conversation.
\end{boxquote}  

Prompt:
\begin{boxquote}{0.9}
You are acting as a teaching assistant for an advanced undergraduate/graduate
algorithms and computer science theory course.

Your role is not to provide full solutions or final answers.

Instead, you should act like we are having a conversation. Ask me a single
question and let me respond. Avoid asking me many questions at once or giving
me a lot of information at one time. Guide me using hints, leading questions,
partial insights, and sanity checks.  Help me debug my own proofs, algorithms,
or reductions. Do not give a complete solution or full proofs. Be patient with
me. Don't give me details so that I can move on to the next part of a problem.
Make sure that I understood before moving on.

The course uses Jeff Erickson's Models of Computation lecture notes, so you should
guide me to answer problems using the definitions, notation, and style from those notes.

Assume I am a serious student trying to learn, not to shortcut the assignment.
I am attaching the assignment, so you should refuse to provide a complete solution
to the problems listed. Of course, you can walk me through the "Solved
Problems" listed at the end if I ask.
\end{boxquote}  



\newpage
%----------------------------------------------------------------------

\headers{CSCI 432/532}{Homework 5 (due February 23)}{Spring 2026}

\begin{enumerate}
%\parindent 1.5em \itemsep 4ex plus 0.5fil

%----------------------------------------------------------------------
\item
Give context-free grammars for the following languages, and clearly explain how
they work and the set of strings generated by each nonterminal. Be sure to
reference the rubric--note that a CFG without an explanation may receive little
or no credit. On the other hand, we do not need a formal proof of correctness.
\begin{enumerate}[(a)]
\item $ \{ \Sym0^a\Sym{10}^b\Sym{10}^c : b = 2a + 2c \}$.
Note that the $^b$ and $^c$ bind only to the \Sym0, not \Sym{10}.
\item $ \{ \Sym0^a\Sym{10}^b\Sym{10}^c : 2b = a + c \}$.
Note that the $^b$ and $^c$ bind only to the \Sym0, not \Sym{10}.
\item The set of all palindromes in $\Sigma^*$ whose lengths are divisible by 7.
\end{enumerate}
\end{enumerate}

\begin{Rubric}{Context-free grammar rubric}
10 points =
\begin{itemize}
\item[+ 2] for a syntactically correct context-free grammar.
\item[+ 4] for a \emph{brief} English explanation of your context-free grammar. This is how you argue that your CFG is correct. We do not want a \emph{transcription}; don't just translate the CFG \emph{notation} into English.
\begin{itemize}
	\item 2 points for an overall explanation and 2 points for giving the language produced by \emph{each} nonterminal.
    \item A CFG without an explanation cannot receive the 4 points for correctness, so its maximum score is 2.
\end{itemize}
\item[+ 4] for correctness
\end{itemize}
\end{Rubric}
% ===========================================================================
\newpage

\subsection*{Solved problems}

Give context-free grammars for the following languages over the alphabet $\Sigma
= \set{\Sym0,\Sym1}$.  Clearly explain how they work and the set of strings
generated by each nonterminal.  Grammars with unclear or missing explanations
may receive little or no credit; on the other hand, we do \emph{not} want formal
proofs of correctness.

\begin{enumerate}\parindent 1.5em
\item
In any string, a \EMPH{run} is a maximal non-empty substring of identical
symbols.  For example, the string $\Sym{0111000011001} = \Sym0^1 \Sym1^3 \Sym0^4
\Sym1^2 \Sym0^2 \Sym1^1$ consists of six runs.

Let $L_a$ be the set of all strings in $\Sigma^*$ that contain two runs of \Sym0s of equal length.  For example, $L_a$ contains the strings \Sym{01101111} and \Sym{01001011100010} (because each of those strings contains more than one run of \Sym0s of length $1$) but $L_a$ does not contain the strings \Sym{000110011011} and \Sym{00000000111}.


\begin{Solution}
\begin{align*}
	S & \to A C B
		&& \text{strings with two blocks of \Sym0s of same length}\\
	A & \to \e \mid X \Sym1
		&&  \text{empty or ends with \Sym1} \\
	B & \to \e \mid \Sym1 X
		&&  \text{empty or starts with \Sym1} \\
	C & \to \Sym0 C \Sym0 \mid \Sym0 D \Sym0
    	&& \text{$\Sym0^n y \Sym0^n$, where $y$ starts and ends with \Sym1}\\
	D & \to \Sym1 \mid \Sym1 X \Sym1
    	&& \text{starts and ends with \Sym1}\\
	X & \to \e \mid \Sym1 X \mid \Sym0 X
		&& \text{all strings:~} (\Sym0+\Sym1)^*
\end{align*}
Every string in $L$ has the form $x\Sym0^n y \Sym0^n z$, where $x$ is either empty or ends with~\Sym1, $y$~starts and ends with \Sym1, and $z$ is either empty or begins with \Sym1.  Nonterminal~$A$ generates the prefix $x$; non-terminal $B$ generates the suffix $z$; nonterminal $C$ generates the matching runs of \Sym0s, and nonterminal $D$ generates the interior string $y$.

The same decomposition can be expressed more compactly as follows:
\begin{align*}
	S & \to B \mid B\Sym1A \mid A\Sym1B \mid A\Sym1B\Sym1A
		&& \text{strings with two blocks of \Sym0s of same length}\\
	A & \to \Sym1A \mid \Sym0A \mid \varepsilon
		&& \text{all strings:~} (\Sym0+\Sym1)^*\\
	B & \to \Sym0B\Sym0 \mid \Sym0\Sym1\Sym0\mid \Sym0\Sym1A\Sym1\Sym0
    	&& \text{$\Sym0^n y \Sym0^n$, where $y$ starts and ends with \Sym1}
\end{align*}
\end{Solution}
\unskip
\begin{rubric}
5 points = 3 for clearly correct grammar + 2 for clear explanation.  These are not the only correct solutions.
\end{rubric}

\item
$L_b = \set{w \in \Sigma^* \mid \text{$w$ is \emph{not} a palindrome}}$.

\begin{Solution}
\begin{align*}
	S & \to \Sym0S\Sym0 \mid \Sym0S\Sym1 \mid \Sym1S\Sym0 \mid \Sym1S\Sym1 \mid A
			&& \text{non-palindromes}\\
	A & \to \Sym0B\Sym1 \mid \Sym1B\Sym0
			&& \text{start and end with different symbols}\\
	B & \to \Sym0B \mid \Sym1B \mid \varepsilon
			&& \text{all strings}
\end{align*}
Every non-palindrome $w$ can be decomposed as either $w = x\Sym0 y \Sym1 z$ or $w = x\Sym1 y \Sym0 z$, for some substrings $x,y,z$ such that $\abs{x}=\abs{z}$.  Non-terminal $S$ generates the prefix $x$ and matching-length suffix $z$; non-terminal $A$ generates the distinct symbols, and non-terminal $B$ generates the interior substring $y$.
\end{Solution}
\unskip
\begin{Solution}
\begin{align*}
	S & \to \Sym0S\Sym0 \mid \Sym1S\Sym1 \mid A
			&& \text{non-palindromes}\\
	A & \to \Sym0B\Sym1 \mid \Sym1B\Sym0
			&& \text{start and end with different symbols}\\
	B & \to \Sym0B \mid \Sym1B \mid \varepsilon
			&& \text{all strings}
\end{align*}
Every non-palindrome $w$ must have a prefix $x$ and a substring $y$ such that either $w = x \Sym0 y \Sym1 x^R$ or $w = x\Sym1 y \Sym0 x^R$.  Specifically, $x$ is the longest common prefix of $w$ and~$w^R$.  In the first case, the grammar generates $w$ as follows:
\[
	S \Tostar x\,A\,x^R \To x \,\Sym0 B\Sym1\, x^R \Tostar x\Sym0y\Sym1x^R = w
\]
The derivation for $w = x\Sym1 y \Sym0 x^R$ is similar. 
\end{Solution}
\unskip
\begin{rubric}
5 points = 3 for clearly correct grammar + 2 for clear explanation.  These are not the only correct solutions.
\end{rubric}
\end{enumerate}




\end{document}